\scanpage{131}
\mspage{130}

1) \deleted{Sectionnage}\inserted{Découpage} (d'arêtes) et recollage (de bouts d'arcs libres)
$\Longleftrightarrow$ normalisation partielle et recollage par identifications)

Si $G$ et $G'$ se correspondent, ($G'$ déduit de $G$ par découpage)
\[
M_G\simeq M_{G'},\qquad \widehat M_G\simeq\widehat M_{G'}\qquad
\mathcal X_{G'}\longrightarrow\mathcal X_G,\quad \widehat{\mathcal X}_{G'}-\widehat{\mathcal X}_G
\]
\uline{NB} On a $S=S'$ et $g_s=g'_s$, $\widehat\nu_s=\widehat\nu'_s$ --- mais \deleted{$\unsure{\nu}$} $\nu_s,\vec\nu_s$ \uline{changent}.

Cette opération sert surtout pour exprimer \inserted{\unsure{un|ou}} $M_G$ -- $\widehat M_G$ comme un produit $\prod_{s\in S}M_{g_s,\unsure{\widehat{\vec A}_s}}$ resp $\prod_s\widehat M_{g_s,\widehat A_s}$.

2) Hom. ordinaire de \inserted{MD-}graphes \deleted{pondérés} (hom de graphes, compatible avec pondération) $\unsure{G'|G}\xrightarrow{u}\unsure{G|G'}$ (\unsure{carré|avec}
\diagram{\begin{tikzcd}
S\arrow[r]\arrow[dr,"\uline{g}"']&S'\arrow[d,"g'"]\\
&\mathbb N
\end{tikzcd}}
commutatif). On a alors
\[
M_G\longrightarrow M_{G'},\qquad
\deleted{\widehat M_G\longrightarrow\widehat M_{G'}}
\]
\[
\mathcal X_G\mathbin{\times}_{M_{G'}}\unsure{M_{G'}|M_G}\longrightarrow\mathcal X_G,\qquad
\deleted{\widehat{\mathcal X}_G\mathbin{\times}_{\widehat M_{G'}}\widehat M_G\longrightarrow\widehat{\mathcal X}_G}
\]
\margin{ne passe pas aux compactifiés}
(auquel d'ailleurs $\unsure{M_{G'}\longrightarrow M_G}$ se réduit)

\margin{Mais il faut supposer que $\forall s'\in S'$, l'application induite $\widehat{\vec A}'_{s'}\longrightarrow\widehat{\vec A}_s$ ($s=u(s')$) soit injective --- On dira que $u$ est une \uline{immersion}}

Cas particuliers intéressants : les hom de transition \deleted{\unsure{fonctorialité}} entre les $M_{g\nu}^{!}$ (qui déjà ne passent pas aux compactifiés !)

Hom d'induction de MD-graphes $G',G$, $u:G'-G$
\diagram{\begin{tikzcd}[column sep=large]
S'\arrow[r,"u_0"]&S
\end{tikzcd}}
\diagram{\begin{tikzcd}[column sep=large]
\widehat{\vec A}'=I'\amalg\vec A'\arrow[r,"u_1"]&\widehat{\vec A}=I\amalg\vec A
\end{tikzcd}}

satisfaisant

a) compatibilité aux applications \uline{origine}
\diagram{\begin{tikzcd}
\widehat{\vec A}'\arrow[r,"u_1"]\arrow[d,"o'"']&\widehat{\vec A}\arrow[d,"o"]\\
S'\arrow[r,"u_0"']&S
\end{tikzcd}}
comm.

(mais pas compatible aux applications extrémité, ni aux $\sigma$\inserted{$\sigma'$})

b) $\overbrace{u_1^{-1}(I)}^{I'_0}\subset I'$ i.e. $u_1^{-1}(\vec A)\supset\vec A'$ i.e. $u_1(\vec A')\subset\vec A$
et on veut que sur $\vec A'$, $u_1$ commute à $\sigma,\sigma'$.

\deleted{c) \unsure{Rq} $I'_1=I'\setminus I'_0=I'\cap u_1^{-1}(\vec A)$, i.e. on veut que $u_1(I'_1)\subset\vec A$ soit stable par $\sigma$, et $u_1|I'_1$ injectif}

\inserted{c) $u_1:\widehat{\vec A}'_{s'}\longrightarrow\widehat{\vec A}_s$ ($s=u_0(s')$) est \uline{injectif}}
